\begin{figure}[ht]
\centering
\begin{tikzpicture}[x=1.4cm, y=0.42cm, font=\footnotesize]

% log-log: x = log10(n), y = log10(max error). n = 10,20,40,80 -> log10 ~ 1..1.9
% Centred slope -2; upwind slope -1.

% Axes
\draw[->, thick] (0.8,-12) -- (3.3,-12) node[right] {$\log_{10} n$};
\draw[->, thick] (0.8,-12) -- (0.8,1) node[above, align=center]
  {$\log_{10}(\text{max error})$};

% x ticks (n = 10, 20, 40, 80, 160)
\foreach \k/\lab in {1/10, 1.301/20, 1.602/40, 1.903/80, 2.204/160} {
  \draw (\k,-12) -- (\k,-12.3);
  \node[anchor=north, font=\scriptsize] at (\k,-12.3) {$\lab$};
}
% y ticks
\foreach \d in {0,-3,-6,-9,-12} {
  \draw (0.8,\d) -- (0.7,\d);
  \node[anchor=east, font=\scriptsize] at (0.7,\d) {$10^{\d}$};
}

% Centred scheme: slope -2. At n=10 (x=1) error 10^{-1.7}; at n=80 (x=1.903) error 10^{-3.5}
% line through (1,-1.7) and (2.204,-4.1) slope = (-4.1+1.7)/(2.204-1)= -2.0
\draw[thick, engblue] (1,-1.7) -- (2.204,-4.1);
\fill[engblue] (1,-1.7) circle (1.6pt);
\fill[engblue] (1.301,-2.3) circle (1.6pt);
\fill[engblue] (1.602,-2.9) circle (1.6pt);
\fill[engblue] (1.903,-3.5) circle (1.6pt);
\node[engblue, font=\scriptsize, anchor=west] at (1.7,-2.7)
  {centred, slope $-2$};

% Upwind scheme: slope -1. through (1,-1.0) and (2.204,-2.2)
\draw[thick, engred] (1,-1.0) -- (2.204,-2.2);
\fill[engred] (1,-1.0) circle (1.6pt);
\fill[engred] (1.301,-1.3) circle (1.6pt);
\fill[engred] (1.602,-1.6) circle (1.6pt);
\fill[engred] (1.903,-1.9) circle (1.6pt);
\node[engred, font=\scriptsize, anchor=west] at (1.7,-1.45)
  {upwind, slope $-1$};

\end{tikzpicture}
\caption[Work-precision diagram for the manufactured solution]{The
work-precision diagram for the convection-diffusion solver on the manufactured
solution $u_{\mathrm{m}}(x) = \sin(\pi x)$ at $P = 1$, on log-log axes. Each
grid halving (rightward by $\log_{10} 2$) drops the centred-scheme error by a
factor of four (slope $-2$) and the upwind-scheme error by a factor of two
(slope $-1$), the machine-checked confirmation of each scheme's order of
accuracy. A line that fails to reach its predicted slope signals an
implementation bug, a boundary-condition error, or an unresolved feature.
Reading the slope off this diagram is the single most useful verification an
engineer runs on a new discretisation. Source: authors' analysis using the
method of manufactured solutions.}
\label{fig:vol02:ch16:mms-convergence}
\end{figure}
